\section{Zero-correlation-length conditions}

\begin{definition}[Single-block local orthogonality convention]%
    \label{def:is_locally_orthogonal}
    \lean{MPSTensor.IsLocallyOrthogonal}
    \leanok
    \uses{def:mps_transfer_idempotence}
    For the single-block convention used here, local orthogonality means
    $\E_A^2=\E_A$. The source local-orthogonality condition for a BNT family is
    instead the collection of mixed-sector equations $\E_{j,j'}=0$ for distinct
    components. Thus this definition records the one-block idempotence
    convention.
\end{definition}

\begin{definition}[BNT local orthogonality]%
    \label{def:is_bnt_locally_orthogonal}
    \lean{MPSTensor.IsBNTLocallyOrthogonal}
    \leanok
    \uses{def:mixed_transfer_rect}
    Let $(A_j)_j$ be the BNT components of a tensor in canonical form. The
    family is \emph{locally orthogonal} when, for every pair of distinct
    components,
    \begin{align}
        \E_{j,j'}
        &=\sum_i A_j^i\otimes\overline{A_{j'}^i}=0.
        \notag
    \end{align}
    This is the mixed-sector condition of
    \cite[Definition~3.5]{Cirac2016MPDO_arXiv}.
\end{definition}

\begin{definition}[Virtual-insertion BNT zero correlation length]%
    \label{def:is_bnt_zcl}
    \lean{MPSTensor.IsBNTZCL}
    \leanok
    \uses{def:bnt_cpsv, def:is_cid, def:is_bnt_locally_orthogonal}
    Let $(A_j)_j$ be a basis of normal tensors for $A$. The BNT
    zero-correlation-length condition is that $\operatorname{CID}(A)$ and, for
    every pair of distinct BNT components, $\E_{j,j'}=0$. This
    virtual-insertion condition is not
    \cite[Definition~3.6]{Cirac2016MPDO_arXiv}, which quantifies over physical
    observables on all disjoint regions.
\end{definition}

\begin{definition}[Positive-gap BNT zero correlation length]%
    \label{def:is_positive_gap_bnt_zcl}
    \lean{MPSTensor.IsPositiveGapBNTZCL}
    \leanok
    \uses{def:bnt_cpsv, def:is_positive_gap_physical_cid,
        def:is_bnt_locally_orthogonal}
    A family $(A_j)_j$ of MPS tensors has positive-gap BNT zero correlation
    length for $A$ when $(A_j)_j$ is a basis of normal tensors for $A$ in the
    sense of Definition~\ref{def:bnt_cpsv}, its components are locally
    orthogonal, and the physical correlations of $A$ are independent of
    separation whenever both complementary gaps are positive. Adjacent regions
    are not included.
\end{definition}

\begin{lemma}[Positive-gap BNT zero correlation length, unfolded]%
    \label{lem:is_positive_gap_bnt_zcl_iff}
    \lean{MPSTensor.isPositiveGapBNTZCL_iff}
    \leanok
    \uses{def:is_positive_gap_bnt_zcl}
    Positive-gap BNT zero correlation length is equivalent to the conjunction
    of the BNT relation, positive-gap physical distance independence, and BNT
    local orthogonality.
\end{lemma}
\begin{proof}\leanok
    This is the defining conjunction.
\end{proof}

\begin{definition}[BNT zero correlation length]%
    \label{def:source_bnt_zcl}
    \lean{MPSTensor.IsPhysicalBNTZCL}
    \leanok
    \uses{def:bnt_cpsv, def:source_physical_cid,
        def:is_bnt_locally_orthogonal}
    A BNT family has zero correlation length when the associated matrix-product
    vectors have correlations independent of distance for all disjoint physical
    regions and the BNT components are locally orthogonal. This is
    \cite[Definition~3.6]{Cirac2016MPDO_arXiv}.
\end{definition}

\begin{theorem}[Unequal raw weights obstruct the ZCL--RFP equivalence]%
    \label{thm:halved_weight_zcl_not_rfp}
    \lean{MPSTensor.halvedWeightTensor_counterexample_to_unrestricted_zcl_iff_rfp}
    \lean{MPSTensor.halvedWeightTensor_isPhysicalBNTZCL}
    \lean{MPSTensor.halvedWeightTensor_eq_halvedDecomp_toTensor}
    \lean{MPSTensor.halvedWeightTensor_mpv}
    \lean{MPSTensor.halvedWeightTensor_sameMPV₂_halvedDecomp}
    \lean{MPSTensor.halvedWeightTensor_not_isTransferIdempotent}
    \leanok
    \uses{def:source_bnt_zcl, def:mps_transfer_idempotence, def:sector_bnt_cf}
    Let $P$ be the canonical decomposition with one scalar basis tensor
    $C^0=(1)$, two one-dimensional copies, and respective weights $1$ and
    $1/2$. Write $A=\mathcal A(P)$ for the tensor represented by $P$. Relative
    to these sector coordinates, $A^0$ has block form
    \begin{align}
        A^0
        &=
        \begin{pmatrix}
            1 & 0 \\
            0 & \tfrac12
        \end{pmatrix}.
        \notag
    \end{align}
    The weights $(1,1/2)$ satisfy the raw-weight BNT canonical-form
    normalization, and the corresponding canonical decomposition represents
    $A$. Its MPV coefficient at length~$N$ is $V^{(N)}(A)=1+2^{-N}$. The BNT
    has one component, so local orthogonality is vacuous. Since the physical
    space has dimension one, all physical correlations are independent of
    separation. Thus $A$ has zero correlation length in the sense of
    Definition~\ref{def:source_bnt_zcl}, but $A$ is not a renormalization fixed
    point. Consequently \cite[Theorem~3.8]{Cirac2016MPDO_arXiv} is false under
    its stated raw-weight BNT normalization.
\end{theorem}
\begin{proof}\leanok
    The $(0,0)$ and $(1,1)$ entries of the one-letter blocking equation
    $(A^0)^2=vA^0$ would give $1=v\cdot 1$ and
    $\frac14=v\cdot\frac12$. The first equality forces $v=1$, while
    substitution in the second gives $1/4=1/2$, a contradiction. Hence no such
    scalar exists. The physical distance-independence assertion follows because
    every observable on a one-letter physical space is scalar.
\end{proof}

\begin{lemma}[Source BNT zero correlation length implies its positive-gap form]%
    \label{lem:is_positive_gap_bnt_zcl_of_source_bnt_zcl}
    \lean{MPSTensor.isPositiveGapBNTZCL_of_isPhysicalBNTZCL}
    \leanok
    \uses{def:source_bnt_zcl, def:is_positive_gap_bnt_zcl}
    If a BNT family has zero correlation length for all disjoint physical
    regions, then it has positive-gap BNT zero correlation length.
\end{lemma}
\begin{proof}\leanok
    \uses{lem:is_positive_gap_physical_cid_of_physical_cid}
    Restrict the physical distance-independence condition to configurations in
    which both complementary gaps are positive. The BNT relation and local
    orthogonality are unchanged.
\end{proof}

\begin{theorem}[Direct-sum RFP implies positive-gap BNT zero correlation length]%
    \label{thm:is_positive_gap_bnt_zcl_of_is_rfp_direct_sum}
    \lean{MPSTensor.isPositiveGapBNTZCL_of_isTransferIdempotent_directSum}
    \leanok
    \uses{def:is_positive_gap_bnt_zcl, def:mps_transfer_idempotence}
    Let $(A_j)_j$ be nonzero-dimensional, irreducible, left-canonical,
    pairwise gauge-phase-distinct BNT components. If the direct-sum tensor
    $A=\bigoplus_j A_j$ satisfies $\E_A^2=\E_A$, then $A$ has positive-gap
    physical CID and $\E_{j,j'}=0$ for all $j\ne j'$. Thus this explicit
    direct-sum representative has positive-gap BNT zero correlation length.
    This does not include the adjacent-region cases in
    \cite[Theorem~3.8]{Cirac2016MPDO_arXiv}.
\end{theorem}
\begin{proof}\leanok
    \uses{thm:is_positive_gap_physical_cid_of_is_rfp,
        thm:bnt_locally_orthogonal_of_rfp_direct_sum}
    Idempotence implies positive-gap physical CID by
    Theorem~\ref{thm:is_positive_gap_physical_cid_of_is_rfp}. On each
    off-diagonal bond block, idempotence of $\E_A$ gives
    $\E_{j,j'}\circ\E_{j,j'}=\E_{j,j'}$. Distinct irreducible left-canonical
    blocks satisfy $\rho(\E_{j,j'})<1$. Since the spectrum of an idempotent is
    contained in $\{0,1\}$, the conditions
    $\E_{j,j'}^2=\E_{j,j'}$ and $\rho(\E_{j,j'})<1$ imply
    $\E_{j,j'}=0$.
\end{proof}

\begin{lemma}[BNT basis of a canonical form is a basis of normal tensors for
    its direct sum]%
    \label{lem:bnt_cf_basis_is_bnt_for_direct_sum}
    \lean{MPSTensor.IsBNTCanonicalForm.isCPSVBasisOfNormalTensors_basisDirectSum}
    \leanok
    \uses{def:bnt_cpsv, def:sector_bnt_cf}
    Let $P$ be a BNT canonical form with distinct basis tensors $(A_j)_j$, and
    let $A=\bigoplus_j A_j$ be their direct sum with one unit-weight copy of
    each sector. Then $(A_j)_j$ is a basis of normal tensors for $A$: every
    $A_j$ is a normal tensor, at every positive system length $N$
    \begin{align}
        V^{(N)}(A)
        &=\sum_j V^{(N)}(A_j),
        \notag
    \end{align}
    and the states $\{|V^{(N)}(A_j)\rangle\}_j$ are linearly independent for
    all sufficiently large $N$.
\end{lemma}
\begin{proof}\leanok
    \uses{cor:sector_bnt_cf_basis_normal, thm:mpv_decomposition}
    Every basis block is irreducible and left-canonical with normalized
    self-overlap converging to $1$, hence a normal tensor by
    Corollary~\ref{cor:sector_bnt_cf_basis_normal}. The direct sum is the
    unit-weight block-diagonal tensor, so its length-$N$ matrix-product vector
    splits as the sum of the block vectors by
    Theorem~\ref{thm:mpv_decomposition}. Eventual linear independence is part
    of the canonical-form data of $P$.
\end{proof}

\begin{theorem}[Corrected forward direction of the main MPS theorem,
    multiplicity-one representative]%
    \label{thm:rfp_implies_positive_gap_bnt_zcl_cf_basis_direct_sum}
    \lean{MPSTensor.IsBNTCanonicalForm.isTransferIdempotent_basisDirectSum_isPositiveGapBNTZCL}
    \leanok
    \uses{def:is_positive_gap_bnt_zcl, def:sector_bnt_cf,
        def:mps_transfer_idempotence}
    Let $P$ be a BNT canonical form with distinct basis tensors $(A_j)_j$, and
    let $A=\bigoplus_j A_j$ be their direct sum with one unit-weight copy of
    each sector. If $A$ is a renormalization fixed point,
    \begin{align}
        \E_A^2&=\E_A,
        \notag
    \end{align}
    then $A$ has positive-gap BNT zero correlation length: its physical
    correlations are independent of the separation whenever both complementary
    gaps are positive, and for all distinct BNT components $j\ne j'$
    \begin{align}
        \E_{j,j'}&=0.
        \notag
    \end{align}
    This is the implication (i)$\implies$(ii) of
    \cite[Theorem~3.8]{Cirac2016MPDO_arXiv} at the explicit multiplicity-one
    unit-weight representative. The raw-weight unrestricted biconditional is
    false by Theorem~\ref{thm:halved_weight_zcl_not_rfp}, and the
    adjacent-region cases of the source CID are excluded by
    Definition~\ref{def:is_positive_gap_bnt_zcl}.
\end{theorem}
\begin{proof}\leanok
    \uses{lem:bnt_cf_basis_is_bnt_for_direct_sum,
        thm:is_positive_gap_bnt_zcl_of_is_rfp_direct_sum}
    By Lemma~\ref{lem:bnt_cf_basis_is_bnt_for_direct_sum}, the distinct basis
    tensors of $P$ are a basis of normal tensors for $A$. The basis blocks are
    irreducible, left-canonical, and pairwise gauge-phase distinct, so
    Theorem~\ref{thm:is_positive_gap_bnt_zcl_of_is_rfp_direct_sum} applies to
    the direct sum.
\end{proof}
